\section{模型评价}

\subsection{模型优点}

\begin{enumerate}
    \item \textbf{多维度综合优化框架}：本文模型突破了传统单目标优化的局限，同时考量建设周期、总成本和环境影响三个关键维度。通过构建帕累托前沿，决策者可以清晰理解不同策略之间的权衡关系，而非被迫在信息不完整的情况下做出选择。\textbf{运力比例分配策略}从数学上证明了最优配置的存在性和唯一性，为工程实践提供了坚实的理论基础。
    
    \item \textbf{现实约束的精细建模}：模型充分考虑了实际运营中的非理想因素，包括太空电梯的\textbf{2\%}年故障率、\textbf{5\%}系绳动力学折减，以及火箭的\textbf{3\%}发射失败率。这些参数源自工程实践和技术预测，使模型输出具有实际指导价值，而非仅停留在理论最优。敏感性分析进一步验证了结论在参数不确定性下的稳健性。
    
    \item \textbf{可扩展的模块化架构}：模型的数学结构具有良好的可扩展性。若未来银河港数量增加、新发射基地投入运营、或出现第三种运输技术（如核热推进穿梭机），只需在现有框架中添加相应模块，重新求解优化问题即可获得更新后的最优配置。这种模块化设计确保了模型的长期适用性。
    
    \item \textbf{量化的可持续性指标}：模型首次为星际物流提供了完整的碳足迹分析框架。通过计算电梯（\textbf{0.1公吨CO$_2$/公吨}）与火箭（\textbf{4.0公吨CO$_2$/公吨}）之间\textbf{40倍}的碳强度差异，决策者可以将环境影响纳入投资回报评估。这对于需要向公众和监管机构报告环境绩效的航天机构尤为重要。
    
    \item \textbf{全生命周期视角}：模型不仅分析建设阶段的物流需求，还评估了殖民地运营阶段的持续供给成本。水资源供应分析（\textbf{547.5亿}vs \textbf{8212.5亿美元/年}）揭示了太空电梯在长期运营中的决定性成本优势，为基础设施投资的战略规划提供了全周期视角。
\end{enumerate}

\subsection{模型局限}

\begin{enumerate}
    \item \textbf{静态运力假设的简化}：模型假设两系统的年运力在整个建设周期内保持恒定。然而，实际情况中：(1) 太空电梯可能经历建设爬坡期，初期运力低于设计值；(2) 技术进步可能在100+年周期中显著提升火箭载荷能力；(3) 系统老化可能导致后期运力下降。更精细的模型应引入时变运力函数$C(t)$，但这将显著增加计算复杂度。
    
    \item \textbf{成本模型的线性假设}：单位运输成本被处理为常数，忽略了规模经济效应（累计发射量增加带来的成本下降）、学习曲线效应（运营经验积累带来的效率提升）、以及资源稀缺效应（大规模需求可能推高原材料价格）。更复杂的非线性成本模型可能产生不同的最优配置。
    
    \item \textbf{离散调度问题的连续化处理}：模型将材料运输建模为连续流，而实际中火箭发射是离散事件，需考虑发射窗口、轨道力学约束和月球表面接收能力。对于详细任务规划，需在本文宏观优化结果的基础上，进一步开发离散事件仿真模型。
    
    \item \textbf{地缘政治因素的缺失}：模型假设10个国际发射基地可以在统一协调下为月球殖民地服务。实际中，国际合作的不确定性、技术出口管制、以及突发地缘冲突都可能限制某些基地的可用性。模型未纳入这些难以量化的政治风险。
    
    \item \textbf{技术颠覆性创新的不可预测性}：100+年的规划时间跨度远超常规技术预测的可靠范围。核聚变推进、反物质引擎、或其他当前无法预见的技术突破可能从根本上改变太空运输的经济学，使本文的量化结论失去参考价值。模型结论应被视为"基于当前技术路径"的参考，而非不可动摇的规划蓝图。
\end{enumerate}

\subsection{模型改进方向}

针对上述局限，后续研究可从以下方向扩展模型：

\begin{enumerate}
    \item 引入随机过程模型，描述运力和成本参数的时变不确定性
    \item 开发多阶段随机规划模型，在每个决策点根据实际情况调整配置
    \item 构建离散事件仿真系统，验证连续模型结论在实际调度中的可行性
    \item 纳入地缘政治情景分析，评估关键基地不可用时的应急策略
\end{enumerate}
